\documentclass[10pt]{article}
\usepackage[margin=1in]{geometry}
\usepackage{graphicx}
\usepackage{booktabs}
\usepackage{amsmath}
\usepackage{multirow}
\title{Reproduction Report: Optimal Transport Enhanced Cross-City Site Recommendation}
\author{Reproduction in OTCRe Repository}
\date{}
\begin{document}
\maketitle

\begin{abstract}
This report reproduces the full experimental workflow of ``Optimal Transport Enhanced Cross-City Site Recommendation'' on all four target cities (Chicago, NYC, Singapore, Tokyo), including MF-BPR, LightGCN, OTC-MF, and OTC-LightGCN. We follow the same evaluation metrics (Recall@20, nDCG@20), 5-core filtering, and cross-city optimal transport setting, and provide direct paper-vs-reproduction comparisons in table and figure form.
\end{abstract}

\section{Introduction}
The target paper proposes OTC, which uses Gromov-Wasserstein optimal transport to project brand/region representations from source cities to a target city and fuse projected inference with the target-city recommender output.

\section{Methodology}
We implement the pipeline in the same structure as the paper:
\begin{enumerate}
\item Train backbone models (MF-BPR, LightGCN) per city.
\item Compute GW transport plans for brand and region embeddings across city pairs.
\item Perform barycentric projection from source to target embeddings.
\item Fuse source-inferred scores with target scores by a weight $\gamma$.
\item Grid-search $\gamma$ and report best Recall@20 and nDCG@20.
\end{enumerate}

\section{Experimental Settings}
We use OpenSiteRec data from the repository, 5-core brand filtering, all-ranking evaluation protocol, and train/val/test split of 70\%/10\%/20\%. Models are trained with Adam and early stopping on validation nDCG@20.

\section{Main Results}
\begin{table*}[t]
\centering
\caption{Paper-reported results vs reproduced results on OpenSiteRec.}
\label{tab:paper_vs_repro}
\small
\begin{tabular}{l|cc|cc|cc|cc}
\toprule
Method & \multicolumn{2}{c|}{Chicago} & \multicolumn{2}{c|}{NYC} & \multicolumn{2}{c|}{Singapore} & \multicolumn{2}{c}{Tokyo}\\
& R@20 & nDCG@20 & R@20 & nDCG@20 & R@20 & nDCG@20 & R@20 & nDCG@20\\
\midrule
MF (Paper) & 0.2494 & 0.1465 & 0.1702 & 0.0917 & 0.4430 & 0.2351 & 0.1323 & 0.0781\\
MF (Repro) & 0.1485 & 0.0919 & 0.1034 & 0.0535 & 0.3678 & 0.1850 & 0.0749 & 0.0425\\
\midrule
OTC-MF (Paper) & 0.2669 & 0.1599 & 0.1766 & 0.0981 & 0.4574 & 0.2503 & 0.1372 & 0.0796\\
OTC-MF (Repro) & 0.1727 & 0.0969 & 0.1110 & 0.0576 & 0.3574 & 0.1791 & 0.0842 & 0.0485\\
\midrule
LightGCN (Paper) & 0.2875 & 0.1902 & 0.2087 & 0.1088 & 0.5013 & 0.2745 & 0.1751 & 0.1068\\
LightGCN (Repro) & 0.1708 & 0.1086 & 0.1268 & 0.0636 & 0.3698 & 0.1798 & 0.1032 & 0.0577\\
\midrule
OTC-LightGCN (Paper) & 0.2977 & 0.2016 & 0.2108 & 0.1109 & 0.5078 & 0.2831 & 0.1780 & 0.1100\\
OTC-LightGCN (Repro) & 0.1918 & 0.1181 & 0.1269 & 0.0631 & 0.3354 & 0.1680 & 0.1073 & 0.0598\\
\midrule
\bottomrule
\end{tabular}
\end{table*}

\begin{figure*}[t]
\centering
\includegraphics[width=0.95\textwidth]{D:/作业/dLproj/OTCRe/OTCRe/data/raw/results/fig_repro/main_results_repro.png}
\caption{Reproduced main results on four cities, matching paper-style grouped comparison across methods.}
\label{fig:main_repro}
\end{figure*}

\begin{figure}[t]
\centering
\includegraphics[width=0.9\linewidth]{D:/作业/dLproj/OTCRe/OTCRe/data/raw/results/fig_repro/improv_recall_repro.png}
\caption{Relative Recall@20 improvement of OTC over its backbone methods in reproduction.}
\label{fig:improv_repro}
\end{figure}

\section{Gamma Sensitivity}
\begin{figure}[t]
\centering
\includegraphics[width=0.9\linewidth]{D:/作业/dLproj/OTCRe/OTCRe/data/raw/results/fig_repro/gamma_lightgcn_chicago.png}
\caption{Gamma sensitivity curve (OTC-LightGCN on Chicago).}
\label{fig:gamma_lgcn_chicago}
\end{figure}

\begin{figure}[t]
\centering
\includegraphics[width=0.9\linewidth]{D:/作业/dLproj/OTCRe/OTCRe/data/raw/results/fig_repro/gamma_vanillamf_chicago.png}
\caption{Gamma sensitivity curve (OTC-MF on Chicago).}
\label{fig:gamma_mf_chicago}
\end{figure}

\section{Discussion}
The reproduction recovers the full multi-city workflow and partially matches the trend of OTC gains in some cities (e.g., Chicago/Tokyo), while deviating on others (notably Singapore under current implementation). This indicates remaining implementation or setting gaps compared with the original authors' internal code path.

\section{Conclusion}
We complete end-to-end reproduction code and reporting artifacts for all cities and all core methods in the paper, with structured comparison to paper-reported numbers.

\end{document}
